\section{An Exposition of Xia's Framework Based on Classification of Projective Binary Groups}

In this section, we will give an exposition of Xia's framework of classifying all $\Holant(\mathcal{F})$ problems with unknown complexity according to the projective binary group generated by $\mathcal{F}.$
This exposition is short and self-contained.
Specifically, we will give a proof for \cref{thm: binary set is finite group}:

\MingjiTheorem*

Recall in \cref{subsec: Induction framework}, we identify $B(f)$ by a subset of $\mathbf{PU}(2)$, and $B(f)$ is closed under multiplication.
Suppose $b\in B(f)$ (we equivalently say the matrix $M(b)\in B(f)$). 
By \cref{def: generating process}, $B(f)$ is closed under conjugation and transpose. 
So $\overline{M(b)}^{\tt T}\in B(f)$.
Let $\overline{b}^{\tt T}$ denote the binary signature in $B(f)$ with matrix $\overline{M(b)}^{\tt T}$.
Notice that $\overline{M(b)}^{\tt T}M(b)=\lambda I_2$ for some $\lambda\in \mathbb
C^{\times}$ since $M(b)\in \mathbf{PU}(2)$.
So $b^{-1}=\overline{b}^{\tt T}$.
Thus $B(f)$ is a group.
In the following we show it is a \emph{finite} group, thus $B(f)$ belongs to a finite list of subgroups of $\mathbf{SO}(3)$, unless $\Holant(\mathcal{F})$ is \#P-hard.
We first reall
\InterpolationBinary*

By \cref{lm: 2 by 2 interpolation}, we may assume that every binary matrix in $B(f)$ has finite projective order, otherwise we can realize $M(g)=\left[\begin{smallmatrix}
    1 & 0\\
    0 & 0
\end{smallmatrix}\right]$ and thus $\Holant(\mathcal{F})$ is \#P-hard by \cref{lm: odd holant with conjugation}.
The following lemma shows $B(f)$ has finite exponent, i.e., there exists a uniform $N\in \N$ such that $b$ has order less than $N$ for every $b\in B(f)$.

\begin{lemma}\label{lm: finite exponent}
    There exists $N\in \N$, such that $b$ has order less than $N$ for every $b\in B(f)$. 
\end{lemma}
\begin{proof}
    Recall that we assume $f$ takes algebraic complex value. 
    Adjoint all possible values of $f$ into $\Q$, and obtain a finite field extension $E/\mathbb{Q}$.
    All possible values of $b$ also lie in $E$ for all $b\in B(f)$, since gadget operations and signature factorizations do not leave $E.$
    Fix an arbitrary $b\in B(f)$.
    By \cref{lm: 2 by 2 interpolation}, $M(b)$ finite projective order $d$ for some $d\in \N$..
    So the eigenvalues $\lambda_1(b)$ and $\lambda_2(b)$ are non-zero, and $\frac{\lambda_2(b)}{\lambda_1(b)}$ is a primitive $d$-th root of unity.

    Now suppose for a contradiction that for every $n\in\N$, there exists $b_n\in B(f)$, such that $b_n$ has order $d_n>n.$ 
    Then $\frac{\lambda_2(b_n)}{\lambda_1(b_n)}$ is a primitive $d_n$-th root of unity.
    However, $\lambda_1(b_n),\lambda_2(b_n)$ are two roots of the characteristic polynomial  of $M(b)$, which is a polynomial in $E[X]$.
    So $\lambda_2(b_n),\lambda_2(b_n)$ lie in a quadratic extension of $E$, as well as $\frac{\lambda_2(b_n)}{\lambda_1(b_n)}$.
    This contradicts $d_n\to +\infty$ as $n\to +\infty.$
\end{proof}
\begin{theorem}[Burnside. See for example, \cite{rowen2008graduate} Theorem 19A.9]\label{thm: Burnside}
   Suppose $F$ is a field of characteristic 0.
   Then every subgroup $G \subseteq \mathrm{GL}(n,F)$ of finite exponent is a finite group.
\end{theorem}

By \cref{thm: Burnside}, $B(f)$ is a finite group, since the representatives of $B(f)$ can be chosen in $\mathbf{SU}(2)\subseteq\mathrm{GL(2,\mathbb{C})}$.
Thus \cref{thm: binary set is finite group} is proved.
Since $\mathbf{SU}(2)/\{\pm I_2\}\cong \mathbf{SO}(3)$, $B(f)$ belongs to the following list: cyclic groups, dihedral groups, or polyhedral groups.